\documentclass[11pt]{article}

\usepackage[letterpaper,margin=1in]{geometry}
\usepackage[T1]{fontenc}
\usepackage[utf8]{inputenc}
\usepackage{lmodern}
\usepackage{microtype}
\usepackage{amsmath,amssymb,mathtools,bm}
\usepackage{enumitem}
\usepackage{xcolor}
\usepackage{booktabs}
\usepackage{array}
\usepackage{hyperref}
\usepackage[most]{tcolorbox}
\usepackage{fancyhdr}
\usepackage{titlesec}
\usepackage{parskip}

\definecolor{qftblue}{HTML}{245A8D}
\definecolor{qftgreen}{HTML}{2F6B4F}
\definecolor{qftred}{HTML}{A23B3B}
\definecolor{qftgray}{HTML}{F4F6F8}

\hypersetup{
  colorlinks=true,
  linkcolor=qftblue,
  urlcolor=qftblue,
  citecolor=qftgreen,
  pdftitle={PHYS 662 Lecture 5: Lagrangians and the Action Principle},
  pdfauthor={Nima Lashkari}
}

\pagestyle{fancy}
\fancyhf{}
\lhead{PHYS 662: Quantum Field Theory I}
\rhead{Lecture 5}
\cfoot{\thepage}
\setlength{\headheight}{14pt}

\titleformat{\section}{\Large\bfseries\color{qftblue}}{\thesection.}{0.6em}{}
\titleformat{\subsection}{\large\bfseries\color{qftgreen}}{\thesubsection}{0.6em}{}
\titleformat{\subsubsection}{\normalsize\bfseries\color{qftred}}{\thesubsubsection}{0.6em}{}

\setlist[itemize]{leftmargin=2em,itemsep=0.25em,topsep=0.4em}
\setlist[enumerate]{leftmargin=2em,itemsep=0.25em,topsep=0.4em}

\newcommand{\dd}{\mathrm{d}}
\newcommand{\R}{\mathbb{R}}
\newcommand{\Lag}{\mathcal{L}}
\newcommand{\cM}{\mathcal{M}}
\newcommand{\SO}{\mathrm{SO}}
\newcommand{\OO}{\mathrm{O}}

\newtcolorbox{keyidea}[1][]{
  enhanced, breakable, colback=qftgray, colframe=qftblue,
  boxrule=0.8pt, arc=1.5mm, title={#1}, fonttitle=\bfseries,
  coltitle=white
}
\newtcolorbox{remarkbox}[1][]{
  enhanced, breakable, colback=white, colframe=qftgreen,
  boxrule=0.7pt, arc=1.5mm, title={#1}, fonttitle=\bfseries,
  coltitle=white
}
\newtcolorbox{warningbox}[1][]{
  enhanced, breakable, colback=white, colframe=qftred,
  boxrule=0.7pt, arc=1.5mm, title={#1}, fonttitle=\bfseries,
  coltitle=white
}

\begin{document}
\pagenumbering{gobble}
\begin{titlepage}
\centering
\vspace*{1.2in}
{\scshape\Large PHYS 662 --- Quantum Field Theory I\par}
\vspace{0.35in}
{\Huge\bfseries\color{qftblue} Lecture 5\par}
\vspace{0.15in}
{\LARGE\bfseries Lagrangians and the Action Principle\par}
\vspace{0.12in}
{\Large\color{qftgreen} From Particle Mechanics to Local Field Theory\par}
\vspace{0.6in}
{\large Nima Lashkari\par}
{\large Fall 2026\par}
\vfill
\begin{minipage}{0.86\textwidth}
\small
This lecture reformulates classical dynamics in terms of an action. After a
brief review of Lorentz covariance, we derive the Euler--Lagrange equations
for finitely many coordinates, identify the role of endpoint terms, and
connect time-translation symmetry to energy conservation. We then pass from
particle trajectories to local fields, define functional derivatives, and
derive the field-theory Euler--Lagrange equations.

\medskip
Course page: \url{https://nimalashkari.com/quantum-field-theory-2026-part-i/}
\end{minipage}
\vspace{0.5in}
\end{titlepage}

\pagenumbering{roman}
\tableofcontents
\newpage
\pagenumbering{arabic}

\section*{Roadmap}
\addcontentsline{toc}{section}{Roadmap}

The action principle packages dynamics, symmetry, conservation laws, and
locality into a single object. The progression of the slides is
\begin{enumerate}
  \item review Minkowski geometry, Lorentz transformations, and covariant
  fields;
  \item formulate stationary action for a mechanical trajectory;
  \item derive the Euler--Lagrange equations and isolate the endpoint term;
  \item introduce canonical momentum, the Hamiltonian, and energy
  conservation;
  \item replace particle coordinates by fields and a mechanical Lagrangian by
  a local Lagrangian density;
  \item define functional derivatives from a regulated set of variables; and
  \item derive the Euler--Lagrange field equations.
\end{enumerate}

Two distinctions will be important throughout. First, a classical solution
makes the action \emph{stationary}; it need not minimize it. Second, locality
means that the Lagrangian density at a point depends on fields and finitely
many derivatives at that same point, even though the action is a functional
of the complete field configuration.

\section{Relativity and covariance: a compact review}

\subsection{Minkowski vectors and the invariant metric}

Let spacetime have one time and \(d\) spatial dimensions, so \(D=d+1\). We use
units \(c=1\), coordinates
\begin{equation}
  x^\mu=(x^0,\bm{x})=(t,\bm{x}),
\end{equation}
and the mostly-minus metric
\begin{equation}
  \eta_{\mu\nu}=\operatorname{diag}(1,-1,\ldots,-1).
  \label{eq:metric}
\end{equation}
Indices are raised and lowered with \(\eta\):
\begin{equation}
  v_\mu=\eta_{\mu\nu}v^\nu,
  \qquad
  \eta^{\mu\rho}\eta_{\rho\nu}=\delta^\mu{}_\nu.
\end{equation}
Thus, if \(v^\mu=(v^0,\bm v)\), then
\begin{equation}
  v_\mu=(v^0,-\bm v),
  \qquad
  v^2:=v_\mu v^\mu=(v^0)^2-|\bm v|^2.
  \label{eq:MinkowskiNorm}
\end{equation}
For two vectors,
\begin{equation}
  v\cdot w:=\eta_{\mu\nu}v^\mu w^\nu
  =v^0w^0-\bm v\cdot\bm w.
\end{equation}
For four-momentum,
\begin{equation}
  p^\mu=(E,\bm p),
  \qquad
  p_\mu=(E,-\bm p),
  \qquad
  p^2=E^2-|\bm p|^2=m^2.
  \label{eq:massshell}
\end{equation}

\begin{warningbox}[Index placement carries a sign]
With the metric in Eq.~\eqref{eq:metric}, the statement
\(p_\mu=(E,\bm p)\) is incorrect: those are the components of \(p^\mu\).
Lowering the index reverses the spatial signs.
\end{warningbox}

\subsection{Lorentz transformations and the Poincare group}

A Lorentz transformation satisfies
\begin{equation}
  \Lambda^T\eta\Lambda=\eta,
  \qquad \Lambda\in\OO(1,d).
  \label{eq:Lorentz}
\end{equation}
It preserves Eq.~\eqref{eq:MinkowskiNorm}. A spatial rotation lies in the
\(\SO(d)\) subgroup. A boost in the \(x^1\) direction can be written
\begin{equation}
  \begin{pmatrix}x'^0\\x'^1\end{pmatrix}
  =
  \begin{pmatrix}
    \cosh\chi&-\sinh\chi\\
    -\sinh\chi&\cosh\chi
  \end{pmatrix}
  \begin{pmatrix}x^0\\x^1\end{pmatrix},
  \quad
  \beta=\tanh\chi,
  \quad
  \gamma=\cosh\chi.
  \label{eq:boost}
\end{equation}
Changing both off-diagonal signs corresponds to \(\chi\mapsto-\chi\); either
choice is valid if the velocity convention is stated.

Rotations and boosts connected continuously to the identity generate
\(\SO^+(1,d)\), not every disconnected component of \(\OO(1,d)\). Adding
translations gives the Poincare group
\begin{equation}
  \operatorname{ISO}(1,d)
  =\R^{1,d}\rtimes\OO(1,d),
  \label{eq:Poincare}
\end{equation}
or \(\R^{1,d}\rtimes\SO^+(1,d)\) when only the identity component is intended.
The product is semidirect because Lorentz transformations act on translation
vectors.

\subsection{Derivatives and covariant fields}

The partial derivative is a covector,
\begin{equation}
  \partial_\mu:=\frac{\partial}{\partial x^\mu},
  \qquad
  \Box:=\partial_\mu\partial^\mu
  =\partial_t^2-\nabla^2.
\end{equation}
For a passive Poincare transformation \(x'=\Lambda x+a\), a scalar obeys
\begin{equation}
  \phi'(x')=\phi(x)
  =\phi\!\left(\Lambda^{-1}(x'-a)\right),
  \label{eq:scalarPassive}
\end{equation}
whereas a contravariant vector obeys
\begin{equation}
  V'^\mu(x')=\Lambda^\mu{}_\nu V^\nu(x).
  \label{eq:vectorPassive}
\end{equation}
More generally, for a multiplet in a representation \(S(\Lambda)\),
\begin{equation}
  \Phi'(x')=S(\Lambda)\Phi(x).
  \label{eq:generalPassive}
\end{equation}
The corresponding active transformation at a fixed coordinate is
\begin{equation}
  (U(\Lambda,a)\Phi)(x)
  =S(\Lambda)\Phi\!\left(\Lambda^{-1}(x-a)\right).
  \label{eq:generalActive}
\end{equation}

\begin{remarkbox}[Covariance versus invariance]
Coordinates and field components generally change. A law is covariant when
the transformed fields satisfy equations of the same form. An action is
invariant when its numerical value is unchanged. These related statements
should not be conflated with the claim that every component is unchanged.
\end{remarkbox}

\section{The action principle in particle mechanics}

\subsection{The action is a functional of a path}

For generalized coordinates \(q^i(t)\), a path connects prescribed endpoints
over \([t_i,t_f]\). The action is
\begin{equation}
  S[q]=\int_{t_i}^{t_f}\dd t\,L(q^i,\dot q^i,t).
  \label{eq:mechanicalAction}
\end{equation}
It is a \emph{functional}: its input is the entire path, and its output is a
number. A neighboring path is
\begin{equation}
  q^i(t)\longmapsto q^i(t)+\varepsilon\eta^i(t),
\end{equation}
where \(\eta^i\) is arbitrary in the interior. For the standard fixed-endpoint
problem,
\begin{equation}
  \eta^i(t_i)=\eta^i(t_f)=0.
  \label{eq:fixedendpoints}
\end{equation}
The physical path obeys
\begin{equation}
  \left.\frac{\dd}{\dd\varepsilon}
  S[q+\varepsilon\eta]\right|_{\varepsilon=0}=0
  \qquad\text{for all allowed }\eta.
  \label{eq:stationary}
\end{equation}

\begin{warningbox}[Stationary, not necessarily least]
Equation~\eqref{eq:stationary} states stationarity. A classical path can be a
local minimum, a local maximum, or a saddle point of the action. The
traditional phrase principle of least action is mathematically too strong.
\end{warningbox}

\subsection{Variation and the Euler--Lagrange equations}

Varying Eq.~\eqref{eq:mechanicalAction} gives
\begin{align}
  \delta S
  &=\int_{t_i}^{t_f}\dd t\,
  \left(
    \frac{\partial L}{\partial q^i}\delta q^i
    +\frac{\partial L}{\partial\dot q^i}\delta\dot q^i
  \right)\notag\\
  &=\int_{t_i}^{t_f}\dd t\,
  \left[
    \frac{\partial L}{\partial q^i}
    -\frac{\dd}{\dd t}
      \left(\frac{\partial L}{\partial\dot q^i}\right)
  \right]\delta q^i
  +\left[
    \frac{\partial L}{\partial\dot q^i}\delta q^i
  \right]_{t_i}^{t_f}.
  \label{eq:variationMechanics}
\end{align}
Repeated configuration indices are summed. Under
Eq.~\eqref{eq:fixedendpoints}, the endpoint term vanishes. Since the interior
variation is arbitrary,
\begin{equation}
  \boxed{
  \frac{\dd}{\dd t}\left(\frac{\partial L}{\partial\dot q^i}\right)
  -\frac{\partial L}{\partial q^i}=0.}
  \label{eq:ELmechanics}
\end{equation}

\subsection{Endpoint terms and well-posed variations}

Adding a total derivative,
\begin{equation}
  L'(q,\dot q,t)=L(q,\dot q,t)+\frac{\dd F(q,t)}{\dd t},
\end{equation}
changes the action only by
\begin{equation}
  S'[q]=S[q]+F(q(t_f),t_f)-F(q(t_i),t_i).
\end{equation}
It therefore leaves Eq.~\eqref{eq:ELmechanics} unchanged for fixed endpoints.
More general boundary terms are allowed only together with boundary
conditions that make the variational problem well posed. For a first-order
Lagrangian, a boundary term depending on \(\dot q\) generally introduces
\(\delta\dot q\) at the endpoint and cannot simply be ignored.

\section{Momentum, energy, and the boundary one-form}

\subsection{Canonical momentum and the Legendre transform}

Define the canonical momentum
\begin{equation}
  p_i:=\frac{\partial L}{\partial\dot q^i}.
  \label{eq:momentum}
\end{equation}
If the Hessian \(\partial^2L/\partial\dot q^i\partial\dot q^j\) is invertible,
the velocities can be expressed in terms of \((q,p,t)\), and the Hamiltonian is
the Legendre transform
\begin{equation}
  H(q,p,t)=p_i\dot q^i-L(q,\dot q,t).
  \label{eq:Hamiltonian}
\end{equation}
Singular Lagrangians require a constrained Hamiltonian analysis; this is the
generic situation in gauge theory.

Along an Euler--Lagrange trajectory,
\begin{align}
  \frac{\dd H}{\dd t}
  &=\dot p_i\dot q^i+p_i\ddot q^i
  -\frac{\partial L}{\partial q^i}\dot q^i
  -\frac{\partial L}{\partial\dot q^i}\ddot q^i
  -\frac{\partial L}{\partial t}\notag\\
  &=-\frac{\partial L}{\partial t}.
  \label{eq:energyidentity}
\end{align}
Hence \(\partial L/\partial t=0\) implies \(\dd H/\dd t=0\). This is the
time-translation case of Noether's theorem.

\subsection{Noether's theorem: a preview}

Suppose an infinitesimal transformation
\(\delta q^i=\varepsilon\Delta q^i\) changes the Lagrangian by a total
derivative,
\begin{equation}
  \delta L=\varepsilon\frac{\dd F}{\dd t}.
\end{equation}
Then, on shell,
\begin{equation}
  Q=p_i\Delta q^i-F,
  \qquad
  \frac{\dd Q}{\dd t}=0.
  \label{eq:NoetherMechanics}
\end{equation}
Thus it is invariance of the action, possibly up to an endpoint term, that
produces a conserved charge.

\subsection{The endpoint term and symplectic geometry}

Without setting endpoint variations to zero, Eq.~\eqref{eq:variationMechanics}
contains
\begin{equation}
  \left[p_i\delta q^i\right]_{t_i}^{t_f}.
\end{equation}
The phase-space one-form
\begin{equation}
  \vartheta=p_i\dd q^i
  \label{eq:symplecticPotential}
\end{equation}
is called the canonical one-form or symplectic potential. Its exterior
derivative determines the symplectic two-form; with our convention,
\begin{equation}
  \omega=-\dd\vartheta=\dd q^i\wedge\dd p_i.
\end{equation}
The boundary contribution is the evaluation of \(\vartheta\) on the endpoint
variation, not itself the symplectic two-form.

\section{From particle coordinates to local fields}

\subsection{Fields and the principle of locality}

A mechanical system has finitely many coordinates \(q^i(t)\). A field theory
has functions \(\phi^a(x)\) on spacetime; the label \(a\) may denote an
internal multiplet, a spacetime tensor index, or both. The local action is
\begin{equation}
  S[\phi]=\int_\cM \dd^D x\,
  \Lag\bigl(\phi^a(x),\partial_\mu\phi^a(x),x\bigr),
  \qquad D=d+1.
  \label{eq:fieldAction}
\end{equation}
Locality means that the integrand at \(x\) depends on the fields and a finite
number of derivatives at \(x\), rather than on independent fields at distant
points. Higher derivatives can be included systematically in an effective
field theory.

For Lorentz invariance, \(\Lag\) should be a Lorentz scalar (more precisely, a
scalar density on a general spacetime), with all indices contracted. For a
Lorentz transformation, \(|\det\Lambda|=1\), so \(\dd^D x\) is invariant.

\begin{keyidea}[Local density, global functional]
The action is global because it integrates over all spacetime. Its
construction is local because each contribution is built from fields and
derivatives evaluated at one point.
\end{keyidea}

\subsection{Dimensional analysis}

In natural units \(\hbar=c=1\), the action is dimensionless. If mass dimension
is denoted by square brackets,
\begin{equation}
  [\dd^D x]=-D,
  \qquad
  [\partial_\mu]=1,
  \qquad
  [\Lag]=D.
  \label{eq:dimensions}
\end{equation}
For a canonically normalized real scalar kinetic term
\begin{equation}
  \Lag_{\rm kin}=\frac12\partial_\mu\phi\,\partial^\mu\phi,
\end{equation}
one finds
\begin{equation}
  [\phi]=\frac{D-2}{2}.
  \label{eq:scalardimension}
\end{equation}
Dimensions depend on the unit convention; without \(\hbar=c=1\), the action
has dimensions of angular momentum rather than zero.

\section{Functional derivatives}

\subsection{From ordinary derivatives to functional derivatives}

For an ordinary function \(f:\R^n\to\R\),
\begin{equation}
  n^i\partial_i f(x)
  =\lim_{\varepsilon\to0}
  \frac{f(x+\varepsilon n)-f(x)}{\varepsilon}.
\end{equation}
A functional \(F[\phi]\) assigns a number to an entire function. Its first
variation defines the functional derivative:
\begin{equation}
  \delta F[\phi]
  =\int \dd^D x\,
  \frac{\delta F}{\delta\phi^a(x)}\,
  \delta\phi^a(x).
  \label{eq:functionalDerivativeDefinition}
\end{equation}
Equivalently,
\begin{equation}
  \frac{\delta\phi^a(x)}{\delta\phi^b(y)}
  =\delta^a{}_b\,\delta^{(D)}(x-y).
  \label{eq:functionalDelta}
\end{equation}

A lattice regulator makes this definition concrete. If cells have volume
\(\Delta V\) and \(\phi_i\approx\phi(x_i)\), then
\begin{equation}
  \delta F=\sum_i\frac{\partial F}{\partial\phi_i}\delta\phi_i
  \longrightarrow
  \int\dd^D x\,
  \frac{\delta F}{\delta\phi(x)}\delta\phi(x),
\end{equation}
so
\begin{equation}
  \frac{\partial F}{\partial\phi_i}
  \simeq \Delta V\,
  \left.\frac{\delta F}{\delta\phi(x)}\right|_{x=x_i}.
  \label{eq:latticeFunctionalDerivative}
\end{equation}
This cell-volume factor accompanies the transition from a Kronecker delta to
a Dirac delta.

\subsection{Examples}

For the linear functional
\begin{equation}
  F[f]=\int\dd^D y\,f(y),
\end{equation}
Eq.~\eqref{eq:functionalDelta} gives
\begin{equation}
  \frac{\delta F}{\delta f(x)}=1.
\end{equation}
For an integral kernel,
\begin{equation}
  F_x[f]=\int\dd^D y\,G(x,y)f(y),
  \qquad
  \frac{\delta F_x}{\delta f(z)}=G(x,z).
\end{equation}
Finally, in one dimension,
\begin{equation}
  F[\phi]=\int\dd x\,(\partial_x\phi)^2
\end{equation}
has variation
\begin{align}
  \delta F
  &=2\int\dd x\,\partial_x\phi\,\partial_x\delta\phi\notag\\
  &=-2\int\dd x\,(\partial_x^2\phi)\delta\phi
  +\left[2\partial_x\phi\,\delta\phi\right]_{\partial\cM}.
\end{align}
For variations that vanish at the boundary,
\begin{equation}
  \frac{\delta F}{\delta\phi(x)}=-2\partial_x^2\phi(x).
  \label{eq:derivativeFunctional}
\end{equation}

\begin{remarkbox}[General versus local functionals]
A generic functional need not possess a convergent polynomial expansion in
discretized field values. When an expansion exists, locality restricts its
kernels to coincident points and their derivatives. A kernel \(G(x,y)\) with
support far from \(x=y\) describes a nonlocal functional.
\end{remarkbox}

\section{Euler--Lagrange equations for fields}

Consider the first-derivative action in Eq.~\eqref{eq:fieldAction}. Its
variation is
\begin{align}
  \delta S
  &=\int_\cM\dd^D x\,
  \left[
    \frac{\partial\Lag}{\partial\phi^a}\delta\phi^a
    +\frac{\partial\Lag}{\partial(\partial_\mu\phi^a)}
      \partial_\mu\delta\phi^a
  \right]\notag\\
  &=\int_\cM\dd^D x\,
  \left[
    \frac{\partial\Lag}{\partial\phi^a}
    -\partial_\mu\left(
      \frac{\partial\Lag}{\partial(\partial_\mu\phi^a)}
    \right)
  \right]\delta\phi^a\notag\\
  &\quad+
  \int_{\partial\cM}\dd\Sigma_\mu\,
  \frac{\partial\Lag}{\partial(\partial_\mu\phi^a)}
  \delta\phi^a.
  \label{eq:variationField}
\end{align}
If variations vanish on \(\partial\cM\), or if suitable boundary conditions
cancel the last term, stationarity gives
\begin{equation}
  \boxed{
  \frac{\partial\Lag}{\partial\phi^a}
  -\partial_\mu\left(
    \frac{\partial\Lag}{\partial(\partial_\mu\phi^a)}
  \right)=0.}
  \label{eq:ELfield}
\end{equation}
Equivalently, the bulk functional derivative is
\begin{equation}
  \frac{\delta S}{\delta\phi^a(x)}
  =
  \frac{\partial\Lag}{\partial\phi^a}
  -\partial_\mu\left(
    \frac{\partial\Lag}{\partial(\partial_\mu\phi^a)}
  \right).
\end{equation}

\subsection{Example: a real scalar field}

For
\begin{equation}
  \Lag
  =\frac12\partial_\mu\phi\,\partial^\mu\phi
  -\frac12m^2\phi^2-V_{\rm int}(\phi),
  \label{eq:scalarLag}
\end{equation}
the derivatives are
\begin{equation}
  \frac{\partial\Lag}{\partial\phi}
  =-m^2\phi-V_{\rm int}'(\phi),
  \qquad
  \frac{\partial\Lag}{\partial(\partial_\mu\phi)}
  =\partial^\mu\phi.
\end{equation}
Equation~\eqref{eq:ELfield} becomes
\begin{equation}
  (\Box+m^2)\phi+V_{\rm int}'(\phi)=0.
  \label{eq:scalarEOM}
\end{equation}
This example displays the complete logic of the lecture: a local scalar
Lagrangian produces a covariant differential equation through stationary
action.

\section{Questions to think about}

\begin{enumerate}
  \item Why does lowering the index of \(p^\mu\) change the signs of its
  spatial components?
  \item Why is the Poincaré group a semidirect rather than a direct product?
  \item Give an example for which a stationary trajectory is not a local
  minimum of the action.
  \item Which boundary condition makes the endpoint term in
  Eq.~\eqref{eq:variationMechanics} vanish?
  \item Why does adding \(\dd F/\dd t\) leave the bulk equations unchanged?
  \item What assumption is needed to solve \(p_i=\partial L/\partial\dot q^i\)
  for the velocities?
  \item How does Eq.~\eqref{eq:energyidentity} encode time-translation
  symmetry?
  \item In what sense is \(S[\phi]\) global while \(\Lag(x)\) is local?
  \item Why does the lattice definition of a functional derivative contain a
  cell-volume factor?
  \item What boundary data are needed in deriving Eq.~\eqref{eq:ELfield}?
\end{enumerate}

\section*{Summary}
\addcontentsline{toc}{section}{Summary}

Minkowski geometry organizes relativistic covariance through an invariant
metric. Lorentz transformations preserve that metric, and their semidirect
product with translations is the Poincare group. Scalars, vectors, and
general multiplets transform through representations of this spacetime
symmetry.

In mechanics, the action is a functional of a path. Requiring its first
variation to vanish for all allowed interior variations gives the
Euler--Lagrange equations. Integration by parts separates the bulk equations
from an endpoint one-form. The Legendre transform introduces canonical
momenta and the Hamiltonian, while continuous variational symmetries lead to
conserved quantities.

In field theory, finitely many coordinates are replaced by functions on
spacetime, and the action is the integral of a local Lagrangian density.
Functional derivatives are the continuum analogues of partial derivatives
with respect to regulated field values. Applying them to a local action
produces the field-theory Euler--Lagrange equations together with a boundary
term that records the required boundary conditions.

\end{document}
